\section{Finite-Dimensional Realization}
\label{sec:transfer}

This section is proves that finite schemes approximate the limiting $b$-series.  Coefficientwise convergence is true, but it is not strong enough to transfer the inverse-majorant condition, which is an infinite nonlinear condition on the inverse series.  We prove a stronger analytic statement.

\subsection{Locally uniform convergence of finite correlations}

Let $H_N$ denote the arcsine-normalized correlation function of the finite scheme $(f_N,g_N)$, in the standard-normal coordinates of Section~\ref{sec:scheme}:
\[
  H_N(t)={\pi\over2}\E[f_N(\calX)g_N(\calY_t)].
\]
Here $\calX,\calY_t\in\R^{2+|D|N}$ are standard Gaussian vectors with $\E[\calX_i\calY_{t,i}]=t$ coordinatewise.

\begin{proposition}[Local uniform convergence]\label{prop:loc-unif}
The functions $H_N$ extend holomorphically to $\D$, satisfy $|H_N(z)|\le\pi/2$ on $\D$, and converge locally uniformly on $\D$ to the limiting correlation $H$.  Consequently, if
\[
  H_N(z)=\sum_{m\ge1}b_{m,N}z^m,
  \qquad
  H(z)=\sum_{m\ge1}b_mz^m,
\]
then $b_{m,N}\to b_m$ for every fixed $m$.
\end{proposition}

\begin{proof}
For fixed $N$, expand $f_N$ and $g_N$ in the orthonormal tensor Hermite basis on $\R^{2+|D|N}$:
\[
  f_N=\sum_\alpha \widehat f_N(\alpha)\He_\alpha,
  \qquad
  g_N=\sum_\alpha \widehat g_N(\alpha)\He_\alpha.
\]
Mehler's formula gives, for real $t\in(-1,1)$,
\[
  {2\over\pi}H_N(t)=\sum_\alpha \widehat f_N(\alpha)\widehat g_N(\alpha)t^{|\alpha|}.
\]
The same series is absolutely convergent for $|z|<1$, since Cauchy-Schwarz gives
\begin{align*}
  \sum_\alpha |\widehat f_N(\alpha)\widehat g_N(\alpha)| |z|^{|\alpha|}
  &\le
  \left(\sum_\alpha |\widehat f_N(\alpha)|^2 |z|^{|\alpha|}\right)^{1/2}
  \left(\sum_\alpha |\widehat g_N(\alpha)|^2 |z|^{|\alpha|}\right)^{1/2}\\
  &\le \|f_N\|_2\|g_N\|_2=1.
\end{align*}
Thus $H_N$ is holomorphic on $\D$ and $|H_N(z)|\le \pi/2$.

We next identify the pointwise limit on the real interval.  Fix $t\in(-1,1)$.  For every degree $d\in D$, the Hermite covariance identity gives
\[
  \E[\He_d(U_{d,j})\He_d(U_{d,j}')]=t^d.
\]
The multivariate central limit theorem yields
\[
  \bigl(P_{d,N},P_{d,N}'\bigr)_{d\in D}
  \Rightarrow
  \bigl(R_d,R_d'\bigr)_{d\in D},
\]
where $(R_d,R_d')$ is a Gaussian pair with correlation $t^d$, and the pairs are independent over distinct $d$.  This convergence is joint with $(Z,Z',X,X')$, which are independent of the Hermite averages.

The limiting discontinuity set of the product of signs is contained in the union of the two threshold hypersurfaces
\begin{align*}
  z+\eta \He_3(x)+\sum_{d\in D}s_dr_d&=0,\\
  z'-\eta \He_3(x')+\sum_{d\in D}\sigma_ds_dr_d'&=0.
\end{align*}
This union has limiting probability zero: after conditioning on all variables except $Z$, the first threshold has probability zero because $Z$ has a density, and similarly for the second threshold using $Z'$.  The bounded continuous mapping theorem therefore gives $H_N(t)\to H(t)$ for every real $t\in(-1,1)$.

The family $\{H_N\}$ is locally bounded on $\D$ and converges pointwise on a set with an accumulation point.  Vitali's theorem gives locally uniform convergence on $\D$ to the unique holomorphic function agreeing with the real limiting correlation.  Coefficient convergence follows from Cauchy's formula.
\end{proof}

\begin{remark}
Proposition~\ref{prop:loc-unif} justifies the limiting coefficient formula as the fixed-degree limit of genuine finite-dimensional schemes.  But it does not by itself imply that the finite inverse series satisfy $\sum_n|A_{n,N}|\gamma^n\le1$.  Tail control for the inverse series requires the inverse-branch stability theorem below.
\end{remark}

\subsection{Analytic stability of inverse branches}

\begin{theorem}[Finite inverse-majorant transfer]\label{thm:finite-transfer}
Let $F_N,F\in\Hol(\D)$, assume $F_N\to F$ locally uniformly on $\D$, and assume $F_N(0)=F(0)=0$.  Suppose that $G$ is a holomorphic inverse branch of $F$ on a neighborhood of $\ol{D(0,R_0)}$, with values in $\D$:
\[
  F(G(\zeta))=\zeta,
  \qquad |\zeta|\le R_0.
\]
Write
\[
  G(\zeta)=\sum_{n\ge1}A_n\zeta^n.
\]
Let $0<\gamma<R_0$, and assume the strict majorant inequality
\[
  \sum_{n\ge1}|A_n|\gamma^n<1.
\]
Then, for all sufficiently large $N$, $F_N$ has a holomorphic inverse branch at the origin,
\[
  G_N(\zeta)=F_N^{-1}(\zeta)=\sum_{n\ge1}A_{n,N}\zeta^n,
\]
defined on a neighborhood of $\gamma\ol\D$, and
\[
  \sum_{n\ge1}|A_{n,N}|\gamma^n<1.
\]
Moreover $A_{n,N}\to A_n$ for every fixed $n$.
\end{theorem}

\begin{proof}
Choose nested radii
\[
  \gamma<\rho<r<R<R_0.
\]
Let
\[
  \Omega=G(D(0,R)).
\]
Since $G$ is a conformal inverse branch on a neighborhood of $\ol{D(0,R)}$, the boundary $\partial\Omega$ is the image of $|\zeta|=R$.  For $t\in\partial\Omega$, we have $|F(t)|=R$.  Therefore, for all $|\zeta|\le r$,
\[
  |F(t)-\zeta|\ge R-r,
  \qquad t\in\partial\Omega.
\]
The compact contour $\partial\Omega$ is contained in $\D$.  By locally uniform convergence, for all large $N$,
\[
  \sup_{t\in\partial\Omega}|F_N(t)-F(t)|<R-r.
\]
Rouche's theorem then implies that, for every $|\zeta|\le r$, the equations $F_N(t)=\zeta$ and $F(t)=\zeta$ have the same number of solutions in $\Omega$, counted with multiplicity.  The equation $F(t)=\zeta$ has exactly one solution in $\Omega$, namely $t=G(\zeta)$.  Hence $F_N(t)=\zeta$ also has exactly one solution in $\Omega$.  In particular, since $F_N(0)=0$, the zero at $0$ is simple, and these solutions define a holomorphic inverse branch $G_N$ on $D(0,r)$, with values in $\Omega$.

The inverse branches converge uniformly on $\ol{D(0,r)}$.  Indeed, by the argument-principle formula,
\[
  G_N(\zeta)= {1\over 2\pi i}\int_{\partial\Omega}
  {tF_N'(t)\over F_N(t)-\zeta}\,dt,
  \qquad |\zeta|\le r.
\]
The denominators are uniformly bounded away from zero on $\partial\Omega\times\ol{D(0,r)}$, and $F_N'\to F'$ uniformly on $\partial\Omega$.  Thus $G_N\to G$ uniformly on $\ol{D(0,r)}$, which implies fixed inverse-coefficient convergence.

It remains to control the tail.  Put
\[
  B=\sup_{t\in\ol\Omega}|t|<\infty.
\]
Since $G_N(D(0,r))\subseteq\Omega$, Cauchy's estimate on the circle $|\zeta|=\rho$ gives
\[
  |A_{n,N}|\le B\rho^{-n},
  \qquad n\ge1,
\]
for all large $N$.  Let
\[
  \delta=1-\sum_{n\ge1}|A_n|\gamma^n>0.
\]
Choose $M$ so large that
\[
  B\sum_{n>M}\left({\gamma\over\rho}\right)^n< {\delta\over3}.
\]
For all sufficiently large $N$, fixed-coefficient convergence gives
\[
  \sum_{n\le M}|A_{n,N}|\gamma^n
  \le \sum_{n\le M}|A_n|\gamma^n+{\delta\over3}.
\]
Therefore
\begin{align*}
  \sum_{n\ge1}|A_{n,N}|\gamma^n
  &\le \sum_{n\le M}|A_n|\gamma^n+{\delta\over3}
  +B\sum_{n>M}\left({\gamma\over\rho}\right)^n\\
  &< \sum_{n\ge1}|A_n|\gamma^n+{2\delta\over3}<1.
\end{align*}
\end{proof}

\begin{corollary}[Finite schemes from a limiting certificate]\label{cor:finite-scheme}
Assume the limiting correlation $H$ has a holomorphic inverse branch
\[
  G(\zeta)=H^{-1}(\zeta)=\sum_{n\ge1}A_n\zeta^n
\]
with values in $\D$ on a neighborhood of $\ol{D(0,R_0)}$.  Let $0<\gamma<R_0$, and suppose
\[
  \sum_{n\ge1}|A_n|\gamma^n<1.
\]
Then, for all sufficiently large lengths $N$, the genuine finite-dimensional scheme $(f_N,g_N)$ satisfies the inverse-majorant condition at radius $\gamma$.  Consequently
\[
  \KG\le {\pi\over2\gamma}.
\]
\end{corollary}

\begin{proof}
Apply Proposition~\ref{prop:loc-unif} and Theorem~\ref{thm:finite-transfer} with $F_N=H_N$ and $F=H$, then invoke Theorem~\ref{thm:krivine-criterion} for the finite scheme after the scaling convention stated at the start of Section~\ref{sec:scheme}.
\end{proof}

\begin{corollary}[Endpoint and margin]\label{cor:margin}
If the limiting inverse branch is holomorphic with values in $\D$ on a neighborhood of $\Gamma\ol\D$ and
\[
  \sum_{n\ge1}|A_n|\Gamma^n\le1,
\]
then every $\gamma<\Gamma$ transfers to all sufficiently high finite-dimensional schemes.  If the inequality at $\Gamma$ is strict and the same analytic neighborhood is available, then $\Gamma$ itself transfers.
\end{corollary}

\begin{proof}
For $\gamma<\Gamma$, the majorant is strict because $A_1\neq0$:
\[
  \sum_{n\ge1}|A_n|\gamma^n<\sum_{n\ge1}|A_n|\Gamma^n\le1.
\]
Apply Corollary~\ref{cor:finite-scheme}.  The strict endpoint case is the same argument with $\gamma=\Gamma$.
\end{proof}
